\documentclass[
  blockchain,
  computerscience,
  chemistry,
  biochemistry,
  biology
]{genesisl1-paper}

% Remove unused field options. The scientific paper remains two-column by default.
% Use \documentclass[onecolumn,...]{genesisl1-paper} only when necessary.

% -----------------------------------------------------------------------------
% Document identity
% -----------------------------------------------------------------------------
\GenesisTitle{Title of the Scientific Paper}
\GenesisSubtitle{Optional subtitle}
\GenesisRunningTitle{Short paper title}
\GenesisVersion{Draft 0.1}
\GenesisStatus{DRAFT}
\GenesisDate{\today}
\GenesisAuthors{Author One, Author Two and Author Three}
\GenesisAffiliations{Affiliation one; affiliation two; ORCID identifiers where applicable}
\GenesisCorrespondence{corresponding.author@example.org}
\GenesisIdentifier{DOI, preprint ID, repository ID, or internal document identifier}
\GenesisRepository{https://example.org/repository}
\GenesisFields{Selected fields}
\GenesisKeywords{keyword one; keyword two; keyword three}
\GenesisLicense{CC BY 4.0}

% Optional exact-artifact metadata. Leave fields empty when irrelevant.
\GenesisDataIdentifier{}
\GenesisCodeIdentifier{}
\GenesisModelIdentifier{}
\GenesisExecutionIdentifier{}
\GenesisContentHash{}
\GenesisIPFSCID{}
\GenesisIPFSGateway{}

\begin{document}

\begin{GenesisFrontMatter}
  \begin{GenesisStructuredAbstract}
    \AbstractBackground{
      In one or two sentences, state the scientific need, the method or approach,
      the measured parameters, and the problem addressed.
    }
    \AbstractMethods{
      State the number of subjects, specimens, datasets, simulations, or software
      systems and name the principal techniques.
    }
    \AbstractResults{
      Use past tense. Report every numerical result required to support the
      conclusion, including uncertainty and exact sample sizes where relevant.
    }
    \AbstractConclusion{
      State the take-home message without introducing a new result or interpretation.
    }
  \end{GenesisStructuredAbstract}
\end{GenesisFrontMatter}

\section{Introduction}
Explain the general scientific need, the present state of knowledge, and the unresolved
gap. Cite prior work constructively and accurately. State the hypothesis or central
question clearly.

\subsection{Aims}
Number multiple aims when useful:
\begin{enumerate}
  \item Determine ...
  \item Compare ...
  \item Evaluate ...
\end{enumerate}

\section{Methods}
Write in past tense and follow the order: subjects or source population, investigated
materials, experimental or computational methods, preprocessing, outcome measures,
and statistics.

\subsection{Subjects, materials, or datasets}
Report inclusion and exclusion criteria, exact sample sizes, versions, and source
identifiers.

\subsection{Experimental or computational procedure}
Describe the procedure at the level required for interpretation and independent
reproduction.

\subsection{Statistical analysis}
State the estimand, statistical tests, one- or two-tailed design, exact \(n\), definitions
of replicates, uncertainty intervals, and decision thresholds.

\begin{GenesisProvenance}
Record repositories, commits, containers, model versions, transaction identifiers,
block heights, content hashes, or IPFS CIDs only when they materially improve
independent inspection. These records do not replace a complete method description.
\end{GenesisProvenance}

\section{Results}
Write in past tense. Follow the order of the figures and cite every figure in numerical
order. Begin with exclusions, dropouts, failed runs, or successful completion counts,
then report quantitative results.

\begin{GenesisEvidence}
State what was directly observed, measured, calculated, reconstructed, or executed.
Support comparative words such as ``higher'', ``lower'', or ``improved'' with values,
uncertainty, and statistical evidence where applicable.
\end{GenesisEvidence}

\section{Discussion}
Begin with the overall meaning of the study. Explain the specific technical or conceptual
novelty, compare the findings neutrally with prior work, and distinguish differences in
methods, populations, materials, parameters, or conditions.

Recommended calibrated phrases include:
\begin{itemize}
  \item ``Our results provide evidence that ...'' for a supported interpretation.
  \item ``This observation points towards ...'' for a preliminary indication.
  \item ``The findings are consistent with ...'' when alternatives remain plausible.
\end{itemize}

\begin{GenesisInterpretation}
State the interpretation separately from direct observation. Avoid presenting an
interpretation as established fact.
\end{GenesisInterpretation}

\begin{GenesisLimitations}
Although encouraging, this study has limitations. State them near the end of the
Discussion and finish with a proportionate positive statement such as:
``Despite these limitations, the consistency of the observed effect supports ...''
\end{GenesisLimitations}

\section{Conclusion}
Summarise the take-home message. Do not introduce a new result or interpretation.
The conclusion should remain consistent with the abstract.

\section*{Data and code availability}
Describe access conditions, repository identifiers, licences, and embargoes. Use
\verb|\PrintGenesisArtifactRecord| when the structured record is useful.

% \PrintGenesisArtifactRecord

\section*{Author contributions}
Use the CRediT taxonomy or another explicit contribution statement.

\section*{Ethics, privacy, rights, and safety}
Provide the declarations required by the field. Do not place personal data, regulated
plaintext, private keys, credentials, or confidential information in public immutable state.

\section*{Competing interests and funding}
State relevant financial, organisational, governance, founder, and commercial roles.

\section*{Acknowledgements}
Keep acknowledgements factual and concise.

\GenesisEndMatterNotice

\bibliographystyle{plainnat}
\bibliography{references}

\end{document}
